\chapter{\IfLanguageName{dutch}{Fase 3 - Custom made application}{Fase 0 - Custom made application}}
\label{ch:fase3}

In dit hoofdstuk doorlopen we het proces van de beheersapplicatie te maken die beantwoord aan de vragen van deze case.

\section{Analyse}

Voor er van start gegaan kan worden, moet eerst duidelijk zijn welke functionaliteiten de PoC allemaal nodig heeft. We doen dit door middel van een analyse, bestaande uit use cases en een EERD. Door deze op te stellen is er een duidelijk overzicht wat gedaan moet worden en kan de CEO waar nodig bijsturen. 

\subsection{Use cases}

\subsubsection{Receptiebon ingeven}

    \vspace{0.5em}
    
    \begin{itemize}[leftmargin=3.5cm, labelsep=0.5cm, nosep]
        \item[\textbf{Primaire actor:}] Productiechef, financieel bediende
        \item[\textbf{Stakeholders:}] /
        \item[\textbf{Precondities:}] De gebruiker is ingelogd
        \item[\textbf{Postcondities:}] De nieuwe receptiebon is aangemaakt
    \end{itemize}
    
    \paragraph{Normaal verloop}
    \begin{enumerate}[nosep]
        \item Klik
        \item klik nog eens
        \item einde
    \end{enumerate}
    
    \paragraph{Alternatieve verlopen}
    \begin{itemize}[nosep]
        \item[1a.] <Omschrijving alternatief pad...>
    \end{itemize}
    
    \paragraph{Domeinspecifieke regels}
    \textbf{DR\_<Naam\_Regel>}: <Omschrijving domeinregel>
    
    \vspace{1cm}

\subsubsection{Historiek bekijken}

\subsubsection{Stock bekijken}

\subsubsection{Stock filteren}

\subsubsection{Stockcorrectie uitvoeren}

\subsubsection{Inloggen}